%============================================================
%  Bd. 40 - Gruppen
%============================================================
\documentclass[main.tex]{subfiles}
\ifSubfilesClassLoaded{\BandSetupStandalone{B40}}{}
\title{Bd. 40 - Gruppen}
\author{Martin Kunze}
\date{}
\setcounter{file}{40}
\hypersetup{pdftitle={Bd. 40 - Gruppen},pdfauthor={Martin Kunze}}
\begin{document}
% Gemeinsame Originalaussagen; nur Band 40 registriert sie.
\newcommand{\ReconstructionStatement}[1]{\csname ReconstructionStatement@#1\endcsname}

\expandafter\newcommand\csname ReconstructionStatement@PowerGroupUnitsAreSingletons\endcsname{%
\FormulaThmDeltaK[Einheiten einer Gruppenpotenzhalbgruppe sind genau Einermengen]{%
  \begin{aligned}[t]
    &\GroupAlg{G}{\star}{e}\dsep X\in\PowNE G\vdash{}\\[-2pt]
    &\MonoidUnit{\PowNE G,\star_{\mathcal P},\{e\}}{X}
      \leftrightarrow X\in\operatorname{Sing}(G)
  \end{aligned}
}{PowerGroupUnitsAreSingletons}{%
  \DeltaRow{Mengen}{G\dsep e\dsep X\dsep g}
  \DeltaRow{Funktionen}{\star\dsep\star_{\mathcal P}}
}
}

\expandafter\newcommand\csname ReconstructionStatement@PowerGroupsSingletonLayer\endcsname{%
\FormulaThmDeltaK[Potenzisomorphismen zwischen Gruppen erhalten die Einermengenschicht]{%
  \begin{aligned}[t]
    &\GroupAlg{G}{\star}{e}\dsep\GroupAlg{H}{\diamond}{u}\dsep{}\\[-2pt]
    &\SemigroupIso{F}{\PowNE G,\star_{\mathcal P}}{\PowNE H,\diamond_{\mathcal P}}
      \vdash{}\\[-2pt]
    &F[\operatorname{Sing}(G)]=\operatorname{Sing}(H)
  \end{aligned}
}{PowerGroupsSingletonLayer}{%
  \DeltaRow{Mengen}{G\dsep H\dsep e\dsep u\dsep X}
  \DeltaRow{Funktionen}{F\dsep R=F\restriction_{\operatorname{Sing}(G)}\dsep\star\dsep\diamond}
}
}

\expandafter\newcommand\csname ReconstructionStatement@PowerGroupsSingletonTransport\endcsname{%
\FormulaThmDeltaK[Gruppenrekonstruktion mit vorgegebenen Einermengenbildern]{%
  \begin{aligned}[t]
    &\GroupAlg{G}{\star}{e}\dsep\GroupAlg{H}{\diamond}{u}\dsep{}\\[-2pt]
    &\SemigroupIso{F}{\PowNE G,\star_{\mathcal P}}{\PowNE H,\diamond_{\mathcal P}}
      \vdash{}\\[-2pt]
    &\exists\varphi\,\left(\begin{aligned}[t]
      &\GroupIso{\varphi}{G,\star,e}{H,\diamond,u}\\[-2pt]
      &\land\forall g\in G\,(F(\{g\})=\{\varphi(g)\})
    \end{aligned}\right)
  \end{aligned}
}{PowerGroupsSingletonTransport}{%
  \DeltaRow{Mengen}{G\dsep H\dsep e\dsep u\dsep g}
  \DeltaRow{Funktionen}{F\dsep\varphi\dsep\star\dsep\diamond}
}
}

\expandafter\newcommand\csname ReconstructionStatement@GroupPowerFullCarrierAbsorption\endcsname{%
\FormulaThmDeltaK[Der volle Gruppenträger absorbiert nichtleere Teilmengen]{%
  \begin{aligned}[t]&\GroupAlg{G}{\star}{e}\dsep Y\in\PowNE G\vdash{}\\&Y\star_{\mathcal P}G=G\land G\star_{\mathcal P}Y=G\end{aligned}
}{GroupPowerFullCarrierAbsorption}{%
  \DeltaRow{Mengen}{G\dsep e\dsep X\dsep Y\dsep x\dsep y\dsep z\dsep a\dsep b}
  \DeltaRow{Funktionen}{\star\dsep\star_{\mathcal P}}
}
}

\expandafter\newcommand\csname ReconstructionStatement@SemigroupIsoUnitStabilityTransport\endcsname{%
\FormulaThmDeltaK[Isomorphismen transportieren Einheitenstabilität]{%
  \begin{aligned}[t]&\MonoidAlg{M}{\star}{e}\dsep\MonoidAlg{N}{\diamond}{d}\dsep{}\\&\SemigroupIso{h}{M,\star}{N,\diamond}\dsep\operatorname{UStab}_{M,\star,e}(z)\vdash{}\\&\operatorname{UStab}_{N,\diamond,d}(h(z))\end{aligned}
}{SemigroupIsoUnitStabilityTransport}{%
  \DeltaRow{Mengen}{M\dsep N\dsep e\dsep d\dsep z\dsep u\dsep w}
  \DeltaRow{Funktionen}{\star\dsep\diamond\dsep h}
}
}

\expandafter\newcommand\csname ReconstructionStatement@PowerMonoidUnitSetStable\endcsname{%
\FormulaThmDeltaK[Die Einheitenmenge ist im Potenzmonoid einheitenstabil]{%
  \MonoidAlg{A}{\star}{e}\vdash \operatorname{UStab}_{\PowNE A,\star_{\mathcal P},\{e\}}(A^\times)
}{PowerMonoidUnitSetStable}{%
  \DeltaRow{Mengen}{A\dsep e\dsep u\dsep W}
  \DeltaRow{Funktionen}{\star\dsep\star_{\mathcal P}}
}
}

\expandafter\newcommand\csname ReconstructionStatement@PowerMonoidUnitStableAbsorption\endcsname{%
\FormulaThmDeltaK[Einheitenstabilität liefert Absorption durch die Einheitenmenge]{%
  \begin{aligned}[t]&\MonoidAlg{A}{\star}{e}\dsep X\in\PowNE A\dsep{}\\&\operatorname{UStab}_{\PowNE A,\star_{\mathcal P},\{e\}}(X)\vdash{}\\&A^\times\star_{\mathcal P}X=X\land X\star_{\mathcal P}A^\times=X\end{aligned}
}{PowerMonoidUnitStableAbsorption}{%
  \DeltaRow{Mengen}{A\dsep e\dsep X\dsep x\dsep z\dsep u}
  \DeltaRow{Funktionen}{\star\dsep\star_{\mathcal P}}

  \DeltaRow{Abkürzungen}{U=A^\times\dsep P=\PowNE A}
}
}

\expandafter\newcommand\csname ReconstructionStatement@PowerMonoidUnitSubsetCriterion\endcsname{%
\FormulaThmDeltaK[Produktkriterium für Teilmengen der Einheitenmenge]{%
  \begin{aligned}[t]&\MonoidAlg{A}{\star}{e}\dsep Y\in\PowNE A\vdash{}\\&Y\in\PowNE{A^\times}\leftrightarrow Y\star_{\mathcal P}A^\times=A^\times\end{aligned}
}{PowerMonoidUnitSubsetCriterion}{%
  \DeltaRow{Mengen}{A\dsep e\dsep Y\dsep y}
  \DeltaRow{Funktionen}{\star\dsep\star_{\mathcal P}}
}
}

\expandafter\newcommand\csname ReconstructionStatement@LargePowerUnitSetImage\endcsname{%
\FormulaThmDeltaK[Bild der gesamten Einheitenmenge]{%
  \begin{aligned}[t]&\MonoidAlg{A}{\star}{e_A}\dsep\MonoidAlg{B}{\diamond}{e_B}\dsep{}\\&\SemigroupIso{F}{\PowNE A,\star_{\mathcal P}}{\PowNE B,\diamond_{\mathcal P}}\vdash F(A^\times)=B^\times\end{aligned}
}{LargePowerUnitSetImage}{%
  \DeltaRow{Mengen}{A\dsep B\dsep e_A\dsep e_B}
  \DeltaRow{Funktionen}{\star\dsep\diamond\dsep\star_{\mathcal P}\dsep\diamond_{\mathcal P}\dsep F}

  \DeltaRow{Abkürzungen}{U=A^\times\dsep V=B^\times}
  \DeltaRow{}{P=\PowNE A\dsep Q=\PowNE B}
}
}

\expandafter\newcommand\csname ReconstructionStatement@LargePowerUnitSubsetImage\endcsname{%
\FormulaThmDeltaK[Bilder nichtleerer Teilmengen der Einheitenmenge]{%
  \begin{aligned}[t]&\MonoidAlg{A}{\star}{e_A}\dsep\MonoidAlg{B}{\diamond}{e_B}\dsep\SemigroupIso{F}{\PowNE A,\star_{\mathcal P}}{\PowNE B,\diamond_{\mathcal P}}\dsep{}\\&F(A^\times)=B^\times\dsep Y\in\PowNE{A^\times}\vdash F(Y)\in\PowNE{B^\times}\end{aligned}
}{LargePowerUnitSubsetImage}{%
  \DeltaRow{Mengen}{A\dsep B\dsep e_A\dsep e_B\dsep Y}
  \DeltaRow{Funktionen}{\star\dsep\diamond\dsep\star_{\mathcal P}\dsep\diamond_{\mathcal P}\dsep F}
}
}

\expandafter\newcommand\csname ReconstructionStatement@LargePowerSemigroupUnitGroupReconstruction\endcsname{%
\FormulaThmDeltaK[Große Potenzhalbgruppen rekonstruieren die
Einheitengruppen]{%
  \begin{aligned}[t]
    &\MonoidAlg{A}{\star}{e_A}\dsep
      \MonoidAlg{B}{\diamond}{e_B}\dsep{}\\[-2pt]
    &\SemigroupIso{F}
      {\PowNE A,\star_{\mathcal P}}
      {\PowNE B,\diamond_{\mathcal P}}\vdash{}\\[-2pt]
    &F(A^\times)=B^\times
      \land F[\PowNE{A^\times}]=\PowNE{B^\times}
  \end{aligned}%
}{LargePowerSemigroupUnitGroupReconstruction}{%
  \DeltaRow{Mengen}{A\dsep B\dsep A^\times\dsep B^\times
    \dsep X\dsep Y\dsep e_A\dsep e_B}
  \DeltaRow{Funktionen}{\star\dsep\diamond\dsep\star_{\mathcal P}
    \dsep\diamond_{\mathcal P}\dsep F}

  \DeltaRow{Abkürzungen}{U=A^\times\dsep V=B^\times}
  \DeltaRow{}{P=\PowNE A\dsep Q=\PowNE B}
  \DeltaRow{}{R=F\restriction_{\PowNE U}}
}
}

\expandafter\newcommand\csname ReconstructionStatement@GroupPowerSemigroupRigidity\endcsname{%
\FormulaThmDeltaK[Einseitige Gruppenstarrheit der großen
Potenzhalbgruppe]{%
  \begin{aligned}[t]
    &\GroupAlg{A}{\star}{e_A}\dsep
      \SemigroupAlg{B}{\diamond}\dsep{}\\[-2pt]
    &\SemigroupIso{F}
      {\PowNE A,\star_{\mathcal P}}
      {\PowNE B,\diamond_{\mathcal P}}\vdash{}\\[-2pt]
    &\exists e_B\in B\,\exists f\,
      \bigl(
        \GroupAlg{B}{\diamond}{e_B}\land
        \GroupIso{f}{A,\star,e_A}{B,\diamond,e_B}
      \bigr)
  \end{aligned}%
}{GroupPowerSemigroupRigidity}{%
  \DeltaRow{Mengen}{A\dsep B\dsep A^\times\dsep B^\times
    \dsep E\dsep X\dsep Y\dsep e_A\dsep e_B}
  \DeltaRow{Funktionen}{\star\dsep\diamond\dsep\star_{\mathcal P}
    \dsep\diamond_{\mathcal P}\dsep F\dsep f}

  \DeltaRow{Abkürzungen}{U=A^\times\dsep V=B^\times}
  \DeltaRow{}{P=\PowNE A\dsep Q=\PowNE B}
}
}

% Der intern benannte Produktzeuge behält seine ursprüngliche Hilfssatznummer.
% Die formale Ableitung steht unverändert in der ausgelagerten Absorptionstabelle.
\newcommand{\ReconstructionProductMemberFormula}{%
  \begin{aligned}[t]
    &\SemigroupAlg{G}{\star}\dsep X,Y\in\PowNE G\dsep x\in X\dsep y\in Y\\[-2pt]
    &\vdash x\star y\in X\star_{\mathcal P}Y
  \end{aligned}%
}
\newcommand{\ReconstructionProductMemberStatement}{%
  \par\Needspace{8\baselineskip}%
  \begingroup
  \setcounter{proofpartnr}{0}%
  \ProofPartSetParent{\theformulaThm}%
  \refstepcounter{proofpartnr}%
  \edef\mxPPLbl{thm:pp:\mxProofPartParent:\arabic{proofpartnr}}%
  \label{\mxPPLbl}%
  \ThmLookupRegisterId{GroupPowerFullCarrierAbsorptionProductMember}{theorem}{\mxPPLbl}%
  \noindent\textbf{Produktzeuge}\hfill\ThmLookupEmitRef{theorem}{\mxPPLbl}\par
  \[\ReconstructionProductMemberFormula\]
  \noindent Die Ableitung bildet den ersten Teil der ausgelagerten Tabelle
  (\ReconstructionProofLink{GroupPowerFullCarrierAbsorption}).\par
  \endgroup
}

\providecommand{\BandXXIXProofRefStack}[1]{%
  \ensuremath{\begin{gathered}[t]#1\end{gathered}}}
\title{Bd. 40 - Gruppen}
\author{Martin Kunze}
\date{}
\hypersetup{pageanchor=false}
\maketitle
\hypersetup{pageanchor=true}
\setcounter{tocdepth}{3}
\tableofcontents
\setcounter{file}{40}
\renewcommand{\FormulaBandID}{40}

\chapter{Einleitung}

Dieser Band ergänzt die Monoide aus Band~38 um beidseitige Inverse.
Er entwickelt die grundlegenden Gruppengesetze, die Eindeutigkeit der
Inversen, Kürzung und Gruppenhomomorphismen. Den Abschluss bilden die
Einheitengruppe eines Monoids und die Rekonstruktion einer Gruppe aus ihrer
großen Potenzhalbgruppe. Abelsche Gruppen werden in Band~41, endliche
Gruppen in Band~42 und Ringe in den Bänden~43--44 behandelt.
Jede Strukturbezeichnung enthält weiterhin die Operation und ihr
neutrales Element, etwa \(\GroupAlg{G}{\circ}{e}\).

\chapter{Gruppen}

\section{Gruppen als Monoide mit Inversen}

\FormulaAxiomDeltaK[Existenz eines beidseitigen Inversen]{%
  \begin{aligned}[t]
    &\MonoidAlg{G}{\circ}{e}\dsep x\in G\vdash{}\\
    &\exists y\in G\,
      \bigl(x\circ y=e\land y\circ x=e\bigr)
  \end{aligned}
}{GroupInverseExistenceLaw}{%
  \DeltaRow{Mengen}{G\dsep e\dsep x\dsep y}
  \DeltaRow{Funktionen}{\circ}
}

\FormulaDefDeltaK[Begriff der Gruppe]{%
  \GroupAlg{G}{\circ}{e}%
}{GroupDef}{%
  \DeltaRow{Mengen}{G\dsep e\dsep x\dsep y}
  \DeltaRow{Funktionen}{\circ}
  \DeltaRow{\textbf{Axiome}}{}
  \DeltaRow{Monoid}{\MonoidAlg{G}{\circ}{e}}
    [\FormulaRefAuto{MonoidDef}[def]]
  \DeltaRow{Beidseitige Inverse}{%
    \begin{aligned}[t]
      &x\in G\vdash{}\\[-2pt]
      &\exists y\in G\,
        \bigl(x\circ y=e\land y\circ x=e\bigr)
    \end{aligned}}
    [\FormulaRefAuto{GroupInverseExistenceLaw}[ax]]
  \DeltaRow{\textbf{Neues Symbol}}{}
  \DeltaPrem{Gruppe}{\GroupAlg{G}{\circ}{e}}
}

\FormulaAxiomDeltaK[Monoidanteil einer Gruppe]{%
  \GroupAlg{G}{\circ}{e}\vdash\MonoidAlg{G}{\circ}{e}%
}{GroupBaseMonoidAxiom}{%
  \DeltaRow{Mengen}{G\dsep e}
  \DeltaRow{Funktionen}{\circ}
}

\FormulaAxiomDeltaK[Inverse in einer Gruppe]{%
  \begin{aligned}[t]
    &\GroupAlg{G}{\circ}{e}\dsep x\in G\vdash{}\\
    &\exists y\in G\,
      \bigl(x\circ y=e\land y\circ x=e\bigr)
  \end{aligned}
}{GroupInverseExistenceAxiom}{%
  \DeltaRow{Mengen}{G\dsep e\dsep x\dsep y}
  \DeltaRow{Funktionen}{\circ}
}

\section{Grundlegende Gruppengesetze}

\FormulaThmDeltaK[Halbgruppenanteil einer Gruppe]{%
  \GroupAlg{G}{\circ}{e}\vdash\SemigroupAlg{G}{\circ}%
}{GroupBaseSemigroup}{%
  \DeltaRow{Mengen}{G\dsep e}
  \DeltaRow{Funktionen}{\circ}
}
\begin{tabproof}
  \proofstep{1}{\GroupAlg{G}{\circ}{e}}{\rA}
  \proofstep{1}{\MonoidAlg{G}{\circ}{e}}{%
    \FormulaRefAuto{GroupBaseMonoidAxiom}[ax]{1}}
  \proofstep{1}{\SemigroupAlg{G}{\circ}}{%
    \FormulaRefAuto{MonoidBaseSemigroupAxiom}[ax]{2}}
\end{tabproof}

\FormulaThmDeltaK[Abgeschlossenheit der Gruppenoperation]{%
  \GroupAlg{G}{\circ}{e}\dsep x,y\in G\vdash x\circ y\in G%
}{GroupClosure}{%
  \DeltaRow{Mengen}{G\dsep e\dsep x\dsep y}
  \DeltaRow{Funktionen}{\circ}
}
\begin{tabproof}
  \proofstep{1}{\GroupAlg{G}{\circ}{e}}{\rA}
  \proofstep{2}{x\in G}{\rA}
  \proofstep{3}{y\in G}{\rA}
  \proofstep{1}{\SemigroupAlg{G}{\circ}}{%
    \FormulaRefAuto{GroupBaseSemigroup}[thm]{1}}
  \proofstep{1,2,3}{x\circ y\in G}{%
    \FormulaRefAuto{SemigroupClosureAxiom}[ax]{4,2,3}}
\end{tabproof}

\FormulaThmDeltaK[Assoziativität der Gruppenoperation]{%
  \begin{aligned}[t]
    &\GroupAlg{G}{\circ}{e}\dsep x,y,z\in G\vdash{}\\
    &(x\circ y)\circ z=x\circ(y\circ z)
  \end{aligned}
}{GroupAssociativity}{%
  \DeltaRow{Mengen}{G\dsep e\dsep x\dsep y\dsep z}
  \DeltaRow{Funktionen}{\circ}
}
\begin{tabproofwide}
  \proofstepwidestar[1]{\GroupAlg{G}{\circ}{e}}{\rA}
  \proofstepwidestar[2]{x\in G}{\rA}
  \proofstepwidestar[3]{y\in G}{\rA}
  \proofstepwidestar[4]{z\in G}{\rA}
  \proofstepwidestar[1]{\SemigroupAlg{G}{\circ}}{%
    \FormulaRefAuto{GroupBaseSemigroup}[thm]{1}}
  \proofstepwidestar[1,2,3,4]{(x\circ y)\circ z=x\circ(y\circ z)}{%
    \FormulaRefAuto{SemigroupAssociativityAxiom}[ax]{5,2,3,4}}
\end{tabproofwide}

\FormulaThmDeltaK[Neutrales Element liegt im Gruppenträger]{%
  \GroupAlg{G}{\circ}{e}\vdash e\in G%
}{GroupIdentityCarrier}{%
  \DeltaRow{Mengen}{G\dsep e}
  \DeltaRow{Funktionen}{\circ}
}
\begin{tabproof}
  \proofstep{1}{\GroupAlg{G}{\circ}{e}}{\rA}
  \proofstep{1}{\MonoidAlg{G}{\circ}{e}}{%
    \FormulaRefAuto{GroupBaseMonoidAxiom}[ax]{1}}
  \proofstep{1}{e\in G}{%
    \FormulaRefAuto{MonoidIdentityCarrierAxiom}[ax]{2}}
\end{tabproof}

\FormulaThmDeltaK[Linksneutralität in einer Gruppe]{%
  \GroupAlg{G}{\circ}{e}\dsep x\in G\vdash e\circ x=x%
}{GroupLeftIdentity}{%
  \DeltaRow{Mengen}{G\dsep e\dsep x}
  \DeltaRow{Funktionen}{\circ}
}
\begin{tabproof}
  \proofstep{1}{\GroupAlg{G}{\circ}{e}}{\rA}
  \proofstep{2}{x\in G}{\rA}
  \proofstep{1}{\MonoidAlg{G}{\circ}{e}}{%
    \FormulaRefAuto{GroupBaseMonoidAxiom}[ax]{1}}
  \proofstep{1,2}{e\circ x=x}{%
    \FormulaRefAuto{MonoidLeftIdentityAxiom}[ax]{3,2}}
\end{tabproof}

\FormulaThmDeltaK[Rechtsneutralität in einer Gruppe]{%
  \GroupAlg{G}{\circ}{e}\dsep x\in G\vdash x\circ e=x%
}{GroupRightIdentity}{%
  \DeltaRow{Mengen}{G\dsep e\dsep x}
  \DeltaRow{Funktionen}{\circ}
}
\begin{tabproof}
  \proofstep{1}{\GroupAlg{G}{\circ}{e}}{\rA}
  \proofstep{2}{x\in G}{\rA}
  \proofstep{1}{\MonoidAlg{G}{\circ}{e}}{%
    \FormulaRefAuto{GroupBaseMonoidAxiom}[ax]{1}}
  \proofstep{1,2}{x\circ e=x}{%
    \FormulaRefAuto{MonoidRightIdentityAxiom}[ax]{3,2}}
\end{tabproof}

\section{Eindeutigkeit des Inversen}

\FormulaThmDeltaK[Eindeutigkeit aus Links- und Rechtsinverse]{%
  \begin{aligned}[t]
    &\GroupAlg{G}{\circ}{e}\dsep x,y,z\in G\dsep
      y\circ x=e\dsep x\circ z=e\\
    &\vdash y=z
  \end{aligned}
}{GroupInverseUnique}{%
  \DeltaRow{Mengen}{G\dsep e\dsep x\dsep y\dsep z}
  \DeltaRow{Funktionen}{\circ}
}
\begin{tabproofsplitwide}
  \proofpartwideindR[Einsetzen der Rechtsinverse]{%
    \begin{aligned}[t]
      &\GroupAlg{G}{\circ}{e}\dsep x,y,z\in G\dsep
        y\circ x=e\dsep x\circ z=e\\
      &\vdash y=y\circ(x\circ z)
    \end{aligned}
  }{%
    \GroupAlg{G}{\circ}{e}\dsep x,y,z\in G\dsep
    y\circ x=e\dsep x\circ z=e
    \vdash y=y\circ(x\circ z)%
  }[GroupInverseUniqueRightInverseSubstitution]
    \proofstepwidestar[1]{\GroupAlg{G}{\circ}{e}}{\rA}
    \proofstepwidestar[2]{x\in G}{\rA}
    \proofstepwidestar[3]{y\in G}{\rA}
    \proofstepwidestar[4]{z\in G}{\rA}
    \proofstepwidestar[5]{y\circ x=e}{\rA}
    \proofstepwidestar[6]{x\circ z=e}{\rA}
    \proofstepwidestar[1,3]{y\circ e=y}{%
      \FormulaRefAuto{GroupRightIdentity}[thm]{1,3}}
    \proofstepwide[1,3]{y}{=}{y\circ e}{%
      \FormulaRefAuto{a=b\vdash b=a}{%
        7}}
    \proofstepwide[6]{y\circ e}{=}{y\circ(x\circ z)}{\rIE{6}}
    \proofstepwidestar[1,2,3,4,5,6]{%
      y=y\circ(x\circ z)}{\rChain{8,9}}
  \closeproofpartwideindR

  \proofpartwideindR[Reduktion mit der Linksinversen]{%
    \begin{aligned}[t]
      &\GroupAlg{G}{\circ}{e}\dsep x,y,z\in G\dsep
        y\circ x=e\dsep x\circ z=e\\
      &\vdash y\circ(x\circ z)=z
    \end{aligned}
  }{%
    \GroupAlg{G}{\circ}{e}\dsep x,y,z\in G\dsep
    y\circ x=e\dsep x\circ z=e
    \vdash y\circ(x\circ z)=z%
  }[GroupInverseUniqueLeftInverseReduction]
    \proofstepwidestar[1]{\GroupAlg{G}{\circ}{e}}{\rA}
    \proofstepwidestar[2]{x\in G}{\rA}
    \proofstepwidestar[3]{y\in G}{\rA}
    \proofstepwidestar[4]{z\in G}{\rA}
    \proofstepwidestar[5]{y\circ x=e}{\rA}
    \proofstepwidestar[6]{x\circ z=e}{\rA}
    \proofstepwidestar[1,2,3,4]{(y\circ x)\circ z=y\circ(x\circ z)}{%
      \FormulaRefAuto{GroupAssociativity}[thm]{1,3,2,4}}
    \proofstepwide[1,2,3,4]{y\circ(x\circ z)}{=}{%
      (y\circ x)\circ z}{%
      \FormulaRefAuto{a=b\vdash b=a}{%
        7}}
    \proofstepwide[5]{(y\circ x)\circ z}{=}{e\circ z}{\rIE{5}}
    \proofstepwide[1,4]{e\circ z}{=}{z}{%
      \FormulaRefAuto{GroupLeftIdentity}[thm]{1,4}}
    \proofstepwidestar[1,2,3,4,5,6]{%
      y\circ(x\circ z)=z}{\rChain{8,9,10}}
  \closeproofpartwideindR

  \proofpartwideindR[Verknüpfung beider Teilrechnungen]{%
    \begin{aligned}[t]
      &\GroupAlg{G}{\circ}{e}\dsep x,y,z\in G\dsep
        y\circ x=e\dsep x\circ z=e\\
      &\vdash y=z
    \end{aligned}
  }{%
    \GroupAlg{G}{\circ}{e}\dsep x,y,z\in G\dsep
    y\circ x=e\dsep x\circ z=e
    \vdash y=z%
  }[GroupInverseUniqueConclusion]
    \proofstepwidestar[1]{\GroupAlg{G}{\circ}{e}}{\rA}
    \proofstepwidestar[2]{x\in G}{\rA}
    \proofstepwidestar[3]{y\in G}{\rA}
    \proofstepwidestar[4]{z\in G}{\rA}
    \proofstepwidestar[5]{y\circ x=e}{\rA}
    \proofstepwidestar[6]{x\circ z=e}{\rA}
    \proofstepwidestar[1,2,3,4,5,6]{%
      y=y\circ(x\circ z)}{%
      \FormulaRefAuto{GroupInverseUniqueRightInverseSubstitution}{%
        1,2,3,4,5,6}}
    \proofstepwidestar[1,2,3,4,5,6]{%
      y\circ(x\circ z)=z}{%
      \FormulaRefAuto{GroupInverseUniqueLeftInverseReduction}{%
        1,2,3,4,5,6}}
    \proofstepwidestar[1,2,3,4,5,6]{y=z}{\rChain{7,8}}
  \closeproofpartwideindR
\end{tabproofsplitwide}

\FormulaDefDeltaK[Inverses Element]{%
  \begin{aligned}[t]
    &\GroupAlg{G}{\circ}{e}\dsep x\in G\vdash{}\\
    &x^{-1}\coloneqq
      \iota y\,
      \bigl(y\in G\land x\circ y=e\land y\circ x=e\bigr)
  \end{aligned}
}{GroupInverseElement}{%
  \DeltaRow{Mengen}{G\dsep e\dsep x\dsep y}
  \DeltaRow{Funktionen}{\circ}
  \DeltaRow{Träger}{x^{-1}\in G}
  \DeltaRow{Rechtsinverse}{x\circ x^{-1}=e}
  \DeltaRow{Linksinverse}{x^{-1}\circ x=e}
  \DeltaRow{Eindeutigkeit}{\FormulaRefAuto{GroupInverseUnique}[thm]}
  \DeltaRow{\textbf{Neues Symbol}}{}
  \DeltaPrem{Inverses Element}{x^{-1}}
}

\FormulaThmDeltaK[Inverses Element liegt im Gruppenträger]{%
  \GroupAlg{G}{\circ}{e}\dsep x\in G\vdash x^{-1}\in G%
}{GroupInverseCarrier}{%
  \DeltaRow{Mengen}{G\dsep e\dsep x}
  \DeltaRow{Funktionen}{\circ}
}
\begin{tabproof}
  \proofstep{1}{\GroupAlg{G}{\circ}{e}}{\rA}
  \proofstep{2}{x\in G}{\rA}
  \proofstep{1,2}{x^{-1}\in G}{%
    \FormulaRefAuto{GroupInverseElement}[def]{1,2}}
\end{tabproof}

\FormulaThmDeltaK[Rechtsinverse Gleichung]{%
  \GroupAlg{G}{\circ}{e}\dsep x\in G\vdash x\circ x^{-1}=e%
}{GroupInverseRight}{%
  \DeltaRow{Mengen}{G\dsep e\dsep x}
  \DeltaRow{Funktionen}{\circ}
}
\begin{tabproof}
  \proofstep{1}{\GroupAlg{G}{\circ}{e}}{\rA}
  \proofstep{2}{x\in G}{\rA}
  \proofstep{1,2}{x\circ x^{-1}=e}{%
    \FormulaRefAuto{GroupInverseElement}[def]{1,2}}
\end{tabproof}

\FormulaThmDeltaK[Linksinverse Gleichung]{%
  \GroupAlg{G}{\circ}{e}\dsep x\in G\vdash x^{-1}\circ x=e%
}{GroupInverseLeft}{%
  \DeltaRow{Mengen}{G\dsep e\dsep x}
  \DeltaRow{Funktionen}{\circ}
}
\begin{tabproof}
  \proofstep{1}{\GroupAlg{G}{\circ}{e}}{\rA}
  \proofstep{2}{x\in G}{\rA}
  \proofstep{1,2}{x^{-1}\circ x=e}{%
    \FormulaRefAuto{GroupInverseElement}[def]{1,2}}
\end{tabproof}

\section{Kürzungsgesetze}

\FormulaThmDeltaK[Linkskürzung in Gruppen]{%
  \begin{aligned}[t]
    &\GroupAlg{G}{\circ}{e}\dsep a,b,c\in G\dsep
      a\circ b=a\circ c\\
    &\vdash b=c
  \end{aligned}
}{GroupLeftCancellation}{%
  \DeltaRow{Mengen}{G\dsep e\dsep a\dsep b\dsep c}
  \DeltaRow{Funktionen}{\circ}
}
\begin{tabproofsplitwide}
  \proofpartwideindR[Linksmultiplikation mit dem Inversen]{%
    \begin{aligned}[t]
      &\GroupAlg{G}{\circ}{e}\dsep a,b,c\in G\dsep
        a\circ b=a\circ c\vdash{}\\[-2pt]
      &(a^{-1}\circ a)\circ b=(a^{-1}\circ a)\circ c
    \end{aligned}
  }{%
    \begin{aligned}[t]
      &\GroupAlg{G}{\circ}{e}\dsep a,b,c\in G\dsep
        a\circ b=a\circ c\vdash{}\\[-2pt]
      &(a^{-1}\circ a)\circ b=(a^{-1}\circ a)\circ c
    \end{aligned}
  }[GroupLeftCancellationInverseMultiplication]
    \proofstepwidestar[1]{\GroupAlg{G}{\circ}{e}}{\rA}
    \proofstepwidestar[2]{a\in G}{\rA}
    \proofstepwidestar[3]{b\in G}{\rA}
    \proofstepwidestar[4]{c\in G}{\rA}
    \proofstepwidestar[5]{a\circ b=a\circ c}{\rA}
    \proofstepwidestar[1,2]{a^{-1}\in G}{%
      \FormulaRefAuto{GroupInverseCarrier}[thm]{1,2}}
    \proofstepwidestar[1,2,5]{%
      a^{-1}\circ(a\circ b)=a^{-1}\circ(a\circ c)}{\rIE{5}}
    \proofstepwide[1,2,3]{%
      (a^{-1}\circ a)\circ b}{=}{a^{-1}\circ(a\circ b)}{%
      \FormulaRefAuto{GroupAssociativity}[thm]{1,6,2,3}}
    \proofstepwidestar[1,2,4]{(a^{-1}\circ a)\circ c=a^{-1}\circ(a\circ c)}{%
      \FormulaRefAuto{GroupAssociativity}[thm]{1,6,2,4}}
    \proofstepwide[1,2,4]{%
      a^{-1}\circ(a\circ c)}{=}{(a^{-1}\circ a)\circ c}{%
      \FormulaRefAuto{a=b\vdash b=a}{%
        9}}
    \proofstepwidestar[1,2,3,4,5]{%
      (a^{-1}\circ a)\circ b=(a^{-1}\circ a)\circ c}{%
      \rChain{8,7,10}}
  \closeproofpartwideindR

  \proofpartwideindR[Reduktion auf das neutrale Element]{%
    \begin{aligned}[t]
      &\GroupAlg{G}{\circ}{e}\dsep a,b,c\in G\dsep
        a\circ b=a\circ c\vdash{}\\[-2pt]
      &e\circ b=e\circ c
    \end{aligned}
  }{%
    \begin{aligned}[t]
      &\GroupAlg{G}{\circ}{e}\dsep a,b,c\in G\dsep
        a\circ b=a\circ c\vdash{}\\[-2pt]
      &e\circ b=e\circ c
    \end{aligned}
  }[GroupLeftCancellationIdentityEquation]
    \proofstepwidestar[1]{\GroupAlg{G}{\circ}{e}}{\rA}
    \proofstepwidestar[2]{a\in G}{\rA}
    \proofstepwidestar[3]{b\in G}{\rA}
    \proofstepwidestar[4]{c\in G}{\rA}
    \proofstepwidestar[5]{a\circ b=a\circ c}{\rA}
    \proofstepwidestar[1,2,3,4,5]{%
      (a^{-1}\circ a)\circ b=(a^{-1}\circ a)\circ c}{%
      \FormulaRefAuto{GroupLeftCancellationInverseMultiplication}[thm]{%
        1,2,3,4,5}}
    \proofstepwidestar[1,2]{a^{-1}\circ a=e}{%
      \FormulaRefAuto{GroupInverseLeft}[thm]{1,2}}
    \proofstepwidestar[1,2]{e=a^{-1}\circ a}{%
      \FormulaRefAuto{a=b\vdash b=a}{7}}
    \proofstepwide[1,2,3]{e\circ b}{=}{(a^{-1}\circ a)\circ b}{\rIE{8}}
    \proofstepwide[1,2,4]{(a^{-1}\circ a)\circ c}{=}{e\circ c}{\rIE{7}}
    \proofstepwidestar[1,2,3,4,5]{e\circ b=e\circ c}{%
      \rChain{9,6,10}}
  \closeproofpartwideindR

  \proofpartwideindR[Entfernen des neutralen Elements]{%
    \begin{aligned}[t]
      &\GroupAlg{G}{\circ}{e}\dsep a,b,c\in G\dsep
        a\circ b=a\circ c\vdash b=c
    \end{aligned}
  }{%
    \begin{aligned}[t]
      &\GroupAlg{G}{\circ}{e}\dsep a,b,c\in G\dsep
        a\circ b=a\circ c\vdash b=c
    \end{aligned}
  }[GroupLeftCancellationConclusion]
    \proofstepwidestar[1]{\GroupAlg{G}{\circ}{e}}{\rA}
    \proofstepwidestar[2]{a\in G}{\rA}
    \proofstepwidestar[3]{b\in G}{\rA}
    \proofstepwidestar[4]{c\in G}{\rA}
    \proofstepwidestar[5]{a\circ b=a\circ c}{\rA}
    \proofstepwidestar[1,2,3,4,5]{e\circ b=e\circ c}{%
      \FormulaRefAuto{GroupLeftCancellationIdentityEquation}[thm]{1,2,3,4,5}}
    \proofstepwidestar[1,3]{e\circ b=b}{%
      \FormulaRefAuto{GroupLeftIdentity}[thm]{1,3}}
    \proofstepwide[1,3]{b}{=}{e\circ b}{%
      \FormulaRefAuto{a=b\vdash b=a}{%
        7}}
    \proofstepwide[1,4]{e\circ c}{=}{c}{%
      \FormulaRefAuto{GroupLeftIdentity}[thm]{1,4}}
    \proofstepwidestar[1,2,3,4,5]{b=c}{\rChain{8,6,9}}
  \closeproofpartwideindR
\end{tabproofsplitwide}

\FormulaThmDeltaK[Rechtskürzung in Gruppen]{%
  \begin{aligned}[t]
    &\GroupAlg{G}{\circ}{e}\dsep a,b,c\in G\dsep
      b\circ a=c\circ a\\
    &\vdash b=c
  \end{aligned}
}{GroupRightCancellation}{%
  \DeltaRow{Mengen}{G\dsep e\dsep a\dsep b\dsep c}
  \DeltaRow{Funktionen}{\circ}
}
\begin{tabproofsplitwide}
  \proofpartwideindR[Rechtsmultiplikation mit dem Inversen]{%
    \begin{aligned}[t]
      &\GroupAlg{G}{\circ}{e}\dsep a,b,c\in G\dsep
        b\circ a=c\circ a\vdash{}\\[-2pt]
      &b\circ(a\circ a^{-1})=c\circ(a\circ a^{-1})
    \end{aligned}
  }{%
    \begin{aligned}[t]
      &\GroupAlg{G}{\circ}{e}\dsep a,b,c\in G\dsep
        b\circ a=c\circ a\vdash{}\\[-2pt]
      &b\circ(a\circ a^{-1})=c\circ(a\circ a^{-1})
    \end{aligned}
  }[GroupRightCancellationInverseMultiplication]
    \proofstepwidestar[1]{\GroupAlg{G}{\circ}{e}}{\rA}
    \proofstepwidestar[2]{a\in G}{\rA}
    \proofstepwidestar[3]{b\in G}{\rA}
    \proofstepwidestar[4]{c\in G}{\rA}
    \proofstepwidestar[5]{b\circ a=c\circ a}{\rA}
    \proofstepwidestar[1,2]{a^{-1}\in G}{%
      \FormulaRefAuto{GroupInverseCarrier}[thm]{1,2}}
    \proofstepwidestar[1,2,5]{%
      (b\circ a)\circ a^{-1}=(c\circ a)\circ a^{-1}}{\rIE{5}}
    \proofstepwidestar[1,2,3]{(b\circ a)\circ a^{-1}=b\circ(a\circ a^{-1})}{%
      \FormulaRefAuto{GroupAssociativity}[thm]{1,3,2,6}}
    \proofstepwide[1,2,3]{%
      b\circ(a\circ a^{-1})}{=}{(b\circ a)\circ a^{-1}}{%
      \FormulaRefAuto{a=b\vdash b=a}{%
        8}}
    \proofstepwide[1,2,4]{%
      (c\circ a)\circ a^{-1}}{=}{c\circ(a\circ a^{-1})}{%
      \FormulaRefAuto{GroupAssociativity}[thm]{1,4,2,6}}
    \proofstepwidestar[1,2,3,4,5]{%
      b\circ(a\circ a^{-1})=c\circ(a\circ a^{-1})}{%
      \rChain{9,7,10}}
  \closeproofpartwideindR

  \proofpartwideindR[Reduktion auf das neutrale Element]{%
    \begin{aligned}[t]
      &\GroupAlg{G}{\circ}{e}\dsep a,b,c\in G\dsep
        b\circ a=c\circ a\vdash{}\\[-2pt]
      &b\circ e=c\circ e
    \end{aligned}
  }{%
    \begin{aligned}[t]
      &\GroupAlg{G}{\circ}{e}\dsep a,b,c\in G\dsep
        b\circ a=c\circ a\vdash{}\\[-2pt]
      &b\circ e=c\circ e
    \end{aligned}
  }[GroupRightCancellationIdentityEquation]
    \proofstepwidestar[1]{\GroupAlg{G}{\circ}{e}}{\rA}
    \proofstepwidestar[2]{a\in G}{\rA}
    \proofstepwidestar[3]{b\in G}{\rA}
    \proofstepwidestar[4]{c\in G}{\rA}
    \proofstepwidestar[5]{b\circ a=c\circ a}{\rA}
    \proofstepwidestar[1,2,3,4,5]{%
      b\circ(a\circ a^{-1})=c\circ(a\circ a^{-1})}{%
      \FormulaRefAuto{GroupRightCancellationInverseMultiplication}[thm]{%
        1,2,3,4,5}}
    \proofstepwidestar[1,2]{a\circ a^{-1}=e}{%
      \FormulaRefAuto{GroupInverseRight}[thm]{1,2}}
    \proofstepwidestar[1,2]{e=a\circ a^{-1}}{%
      \FormulaRefAuto{a=b\vdash b=a}{7}}
    \proofstepwide[1,2,3]{b\circ e}{=}{b\circ(a\circ a^{-1})}{\rIE{8}}
    \proofstepwide[1,2,4]{c\circ(a\circ a^{-1})}{=}{c\circ e}{\rIE{7}}
    \proofstepwidestar[1,2,3,4,5]{b\circ e=c\circ e}{%
      \rChain{9,6,10}}
  \closeproofpartwideindR

  \proofpartwideindR[Entfernen des neutralen Elements]{%
    \begin{aligned}[t]
      &\GroupAlg{G}{\circ}{e}\dsep a,b,c\in G\dsep
        b\circ a=c\circ a\vdash b=c
    \end{aligned}
  }{%
    \begin{aligned}[t]
      &\GroupAlg{G}{\circ}{e}\dsep a,b,c\in G\dsep
        b\circ a=c\circ a\vdash b=c
    \end{aligned}
  }[GroupRightCancellationConclusion]
    \proofstepwidestar[1]{\GroupAlg{G}{\circ}{e}}{\rA}
    \proofstepwidestar[2]{a\in G}{\rA}
    \proofstepwidestar[3]{b\in G}{\rA}
    \proofstepwidestar[4]{c\in G}{\rA}
    \proofstepwidestar[5]{b\circ a=c\circ a}{\rA}
    \proofstepwidestar[1,2,3,4,5]{b\circ e=c\circ e}{%
      \FormulaRefAuto{GroupRightCancellationIdentityEquation}[thm]{1,2,3,4,5}}
    \proofstepwidestar[1,3]{b\circ e=b}{%
      \FormulaRefAuto{GroupRightIdentity}[thm]{1,3}}
    \proofstepwide[1,3]{b}{=}{b\circ e}{%
      \FormulaRefAuto{a=b\vdash b=a}{%
        7}}
    \proofstepwide[1,4]{c\circ e}{=}{c}{%
      \FormulaRefAuto{GroupRightIdentity}[thm]{1,4}}
    \proofstepwidestar[1,2,3,4,5]{b=c}{\rChain{8,6,9}}
  \closeproofpartwideindR
\end{tabproofsplitwide}

\FormulaThmDeltaK[Nichtleere endliche kürzbare Halbgruppen sind Gruppen]{%
  \begin{aligned}[t]
    &\CancellativeSemigroupAlg{A}{\star}\dsep
      \Finite A\dsep A\neq\varnothing\vdash{}\\[-2pt]
    &\exists e\in A\,\GroupAlg{A}{\star}{e}
  \end{aligned}%
}{FiniteCancellativeSemigroupIsGroup}{%
  \DeltaRow{Mengen}{A\dsep a\dsep b\dsep c\dsep d\dsep e\dsep f
    \dsep x\dsep y}
  \DeltaRow{Funktionen}{\star\dsep L_a\dsep R_a}
}
\begingroup
\ProofWideReasonsNearFormula
\begin{tabproofwide}
  \proofstepwidestarbreakkeep[1]{\CancellativeSemigroupAlg{A}{\star}}{%
    \rA}
  \proofstepwidestarbreakkeep[2]{\Finite A}{%
    \rA}
  \proofstepwidestarbreakkeep[3]{A\neq\varnothing}{%
    \rA}
  \proofstepwidestarbreakkeep[1]{\SemigroupAlg{A}{\star}}{%
    \FormulaRefAuto{CancellativeSemigroupBaseAxiom}[ax]{1}}
  \proofstepwidestarbreakkeep[1,2]{\FiniteSemigroupAlg{A}{\star}}{%
    \FormulaRefAuto{FiniteSemigroupDef}[def]{4,2}}
  \proofstepwidestarbreakkeep[1,2,3]{\exists e\in A\,\GroupAlg{A}{\star}{e}}{%
    \FormulaRefAuto{FiniteCancellativeSemigroupGroupExistence}[thm]{5,1,3}}
\end{tabproofwide}
\endgroup


\section{Inverse eines Produkts}

\FormulaThmDeltaK[Inverse eines Produkts]{%
  \begin{aligned}[t]
    &\GroupAlg{G}{\circ}{e}\dsep x,y\in G\vdash{}\\
    &(x\circ y)^{-1}=y^{-1}\circ x^{-1}
  \end{aligned}
}{GroupProductInverse}{%
  \DeltaRow{Mengen}{G\dsep e\dsep x\dsep y}
  \DeltaRow{Funktionen}{\circ}
}
\begin{tabproofsplitwide}
  \proofpartwideindR[Rechtsinverse des Produkts]{%
    \begin{aligned}[t]
      &\GroupAlg{G}{\circ}{e}\dsep x,y\in G\vdash{}\\
      &(x\circ y)\circ(y^{-1}\circ x^{-1})=e
    \end{aligned}
  }{%
    \GroupAlg{G}{\circ}{e}\dsep x,y\in G
    \vdash (x\circ y)\circ(y^{-1}\circ x^{-1})=e%
  }[GroupProductInverseRightCandidate]
    \proofstepwidestarbreakkeep[1]{\GroupAlg{G}{\circ}{e}}{%
      \rA}
    \proofstepwidestarbreakkeep[2]{x\in G}{%
      \rA}
    \proofstepwidestarbreakkeep[3]{y\in G}{%
      \rA}
    \proofstepwidestarbreakkeep[1,2]{x^{-1}\in G}{%
      \FormulaRefAuto{GroupInverseCarrier}[thm]{1,2}}
    \proofstepwidestarbreakkeep[1,3]{y^{-1}\in G}{%
      \FormulaRefAuto{GroupInverseCarrier}[thm]{1,3}}
    \proofstepwidestarbreakkeep[1,2,3]{y^{-1}\circ x^{-1}\in G}{%
      \FormulaRefAuto{GroupClosure}[thm]{1,5,4}}
    \proofstepwidestarbreakkeep[1,2,3]{(x\circ y)\circ(y^{-1}\circ x^{-1})=x\circ(y\circ(y^{-1}\circ x^{-1}))}{%
      \FormulaRefAuto{GroupAssociativity}[thm]{1,2,3,6}}
    \proofstepwidestarbreakkeep[1,2,3]{(y\circ y^{-1})\circ x^{-1}=y\circ(y^{-1}\circ x^{-1})}{%
      \FormulaRefAuto{GroupAssociativity}[thm]{1,3,5,4}}
    \proofstepwidestarbreakkeep[1,3]{y\circ y^{-1}=e}{%
      \FormulaRefAuto{GroupInverseRight}[thm]{1,3}}
    \proofstepwidestarbreakkeep[1,2]{e\circ x^{-1}=x^{-1}}{%
      \FormulaRefAuto{GroupLeftIdentity}[thm]{1,4}}
    \proofstepwidestarbreakkeep[1,2,3]{y\circ(y^{-1}\circ x^{-1})=x^{-1}}{%
      \FormulaRefAuto{a=b\dsep a=c\vdash b=c}[thm]{8,\rIE{9,10}}}
    \proofstepwidestarbreakkeep[1,2,3]{x\circ(y\circ(y^{-1}\circ x^{-1}))=x\circ x^{-1}}{%
      \rIE{11}}
    \proofstepwidestarbreakkeep[1,2]{x\circ x^{-1}=e}{%
      \FormulaRefAuto{GroupInverseRight}[thm]{1,2}}
    \proofstepwidestarbreakkeep[1,2,3]{(x\circ y)\circ(y^{-1}\circ x^{-1})=e}{%
      \rChain{7,12,13}}
  \closeproofpartwideindR

  \proofpartwideindR[Linksinverse des Produkts]{%
    \begin{aligned}[t]
      &\GroupAlg{G}{\circ}{e}\dsep x,y\in G\vdash{}\\
      &(y^{-1}\circ x^{-1})\circ(x\circ y)=e
    \end{aligned}
  }{%
    \GroupAlg{G}{\circ}{e}\dsep x,y\in G
    \vdash (y^{-1}\circ x^{-1})\circ(x\circ y)=e%
  }[GroupProductInverseLeftCandidate]
    \proofstepwidestarbreakkeep[1]{\GroupAlg{G}{\circ}{e}}{%
      \rA}
    \proofstepwidestarbreakkeep[2]{x\in G}{%
      \rA}
    \proofstepwidestarbreakkeep[3]{y\in G}{%
      \rA}
    \proofstepwidestarbreakkeep[1,2]{x^{-1}\in G}{%
      \FormulaRefAuto{GroupInverseCarrier}[thm]{1,2}}
    \proofstepwidestarbreakkeep[1,3]{y^{-1}\in G}{%
      \FormulaRefAuto{GroupInverseCarrier}[thm]{1,3}}
    \proofstepwidestarbreakkeep[1,2,3]{x\circ y\in G}{%
      \FormulaRefAuto{GroupClosure}[thm]{1,2,3}}
    \proofstepwidestarbreakkeep[1,2,3]{(y^{-1}\circ x^{-1})\circ(x\circ y)=y^{-1}\circ(x^{-1}\circ(x\circ y))}{%
      \FormulaRefAuto{GroupAssociativity}[thm]{1,5,4,6}}
    \proofstepwidestarbreakkeep[1,2,3]{(x^{-1}\circ x)\circ y=x^{-1}\circ(x\circ y)}{%
      \FormulaRefAuto{GroupAssociativity}[thm]{1,4,2,3}}
    \proofstepwidestarbreakkeep[1,2]{x^{-1}\circ x=e}{%
      \FormulaRefAuto{GroupInverseLeft}[thm]{1,2}}
    \proofstepwidestarbreakkeep[1,3]{e\circ y=y}{%
      \FormulaRefAuto{GroupLeftIdentity}[thm]{1,3}}
    \proofstepwidestarbreakkeep[1,2,3]{x^{-1}\circ(x\circ y)=y}{%
      \FormulaRefAuto{a=b\dsep a=c\vdash b=c}[thm]{8,\rIE{9,10}}}
    \proofstepwidestarbreakkeep[1,2,3]{y^{-1}\circ(x^{-1}\circ(x\circ y))=y^{-1}\circ y}{%
      \rIE{11}}
    \proofstepwidestarbreakkeep[1,3]{y^{-1}\circ y=e}{%
      \FormulaRefAuto{GroupInverseLeft}[thm]{1,3}}
    \proofstepwidestarbreakkeep[1,2,3]{(y^{-1}\circ x^{-1})\circ(x\circ y)=e}{%
      \rChain{7,12,13}}
  \closeproofpartwideindR

  \proofpartwide{Eindeutigkeit}
    \proofstepwidestar[1]{\GroupAlg{G}{\circ}{e}}{\rA}
    \proofstepwidestar[2]{x\in G}{\rA}
    \proofstepwidestar[3]{y\in G}{\rA}
    \proofstepwidestar[1,2]{x^{-1}\in G}{%
      \FormulaRefAuto{GroupInverseCarrier}[thm]{1,2}}
    \proofstepwidestar[1,3]{y^{-1}\in G}{%
      \FormulaRefAuto{GroupInverseCarrier}[thm]{1,3}}
    \proofstepwidestar[1,2,3]{x\circ y\in G}{%
      \FormulaRefAuto{GroupClosure}[thm]{1,2,3}}
    \proofstepwidestar[1,2,3]{y^{-1}\circ x^{-1}\in G}{%
      \FormulaRefAuto{GroupClosure}[thm]{1,5,4}}
    \proofstepwidestar[1,2,3]{(x\circ y)^{-1}\in G}{%
      \FormulaRefAuto{GroupInverseCarrier}[thm]{1,6}}
    \proofstepwidestar[1,2,3]{%
      (y^{-1}\circ x^{-1})\circ(x\circ y)=e}{%
      \FormulaRefAuto{GroupProductInverseLeftCandidate}[thm]{1,2,3}}
    \proofstepwidestar[1,2,3]{%
      (x\circ y)\circ(x\circ y)^{-1}=e}{%
      \FormulaRefAuto{GroupInverseRight}[thm]{1,6}}
    \proofstepwidestar[1,2,3]{%
      y^{-1}\circ x^{-1}=(x\circ y)^{-1}}{%
      \FormulaRefAuto{GroupInverseUnique}[thm]{1,6,7,8,9,10}}
    \proofstepwidestar[1,2,3]{%
      (x\circ y)^{-1}=y^{-1}\circ x^{-1}}{%
      \FormulaRefAuto{a=b\vdash b=a}{11}}
  \closeproofpartwide
\end{tabproofsplitwide}

\chapter{Homomorphismen von Gruppen}

\section{Gruppenhomomorphismen}

Bei Gruppen genügt die Erhaltung der Gruppenoperation. Anders als bei
beliebigen Monoiden folgen die Erhaltung des neutralen Elements und der
Inversen aus der Existenz von Inversen im Ziel.

\FormulaDefDeltaK[Begriff des Gruppenhomomorphismus]{%
  \GroupHom{f}{G,\circ,e}{H,\diamond,u}%
}{GroupHomDef}{%
  \DeltaRow{Mengen}{G\dsep H\dsep e\dsep u\dsep x\dsep y}
  \DeltaRow{Funktionen}{f\dsep\circ\dsep\diamond}
  \DeltaRow{\textbf{Axiome}}{}
  \DeltaPrem{Quellgruppe}{\GroupAlg{G}{\circ}{e}}
  \DeltaPrem{Zielgruppe}{\GroupAlg{H}{\diamond}{u}}
  \DeltaPrem{Halbgruppenhomomorphismus}{%
    \SemigroupHom{f}{G,\circ}{H,\diamond}}
  \DeltaRow{\textbf{Neues Symbol}}{}
  \DeltaPrem{Gruppenhomomorphismus}{%
    \GroupHom{f}{G,\circ,e}{H,\diamond,u}}
}

\FormulaAxiomDeltaK[Trägerabbildung eines Gruppenhomomorphismus]{%
  \GroupHom{f}{G,\circ,e}{H,\diamond,u}
  \vdash f\colon G\to H%
}{GroupHomFunctionAxiom}{%
  \DeltaRow{Mengen}{G\dsep H\dsep e\dsep u}
  \DeltaRow{Funktionen}{f\dsep\circ\dsep\diamond}
}
\par\smallskip

\FormulaAxiomDeltaK[Operationserhaltung eines Gruppenhomomorphismus]{%
  \GroupHom{f}{G,\circ,e}{H,\diamond,u}\dsep x,y\in G
  \vdash f(x\circ y)=f(x)\diamond f(y)%
}{GroupHomOperationAxiom}{%
  \DeltaRow{Mengen}{G\dsep H\dsep e\dsep u\dsep x\dsep y}
  \DeltaRow{Funktionen}{f\dsep\circ\dsep\diamond}
}

\FormulaThmDeltaK[Gruppenhomomorphismen erhalten das neutrale Element]{%
  \GroupHom{f}{G,\circ,e}{H,\diamond,u}
  \vdash f(e)=u%
}{GroupHomIdentityPreservation}{%
  \DeltaRow{Mengen}{G\dsep H\dsep e\dsep u}
  \DeltaRow{Funktionen}{f\dsep\circ\dsep\diamond}
}
\begin{tabproofwide}
  \proofstepwidestar[1]{%
    \GroupHom{f}{G,\circ,e}{H,\diamond,u}}{\rA}
  \proofstepwidestar[1]{\GroupAlg{G}{\circ}{e}}{%
    \FormulaRefAuto{GroupHomDef}[def]{1}}
  \proofstepwidestar[1]{\GroupAlg{H}{\diamond}{u}}{%
    \FormulaRefAuto{GroupHomDef}[def]{1}}
  \proofstepwidestar[1]{f\colon G\to H}{%
    \FormulaRefAuto{GroupHomFunctionAxiom}[ax]{1}}
  \proofstepwidestar[1]{e\in G}{%
    \FormulaRefAuto{GroupIdentityCarrier}[thm]{2}}
  \proofstepwidestar[1]{f(e)\in H}{%
    \FormulaRefAuto{F\colon A\to B\dsep a\in A\vdash F(a)\in B}{4,5}}
  \proofstepwidestar[1]{e\circ e=e}{%
    \FormulaRefAuto{GroupLeftIdentity}[thm]{2,5}}
  \proofstepwidestar[1]{f(e\circ e)=f(e)}{\rIE{7}}
  \proofstepwidestar[1]{f(e\circ e)=f(e)\diamond f(e)}{%
    \FormulaRefAuto{GroupHomOperationAxiom}[ax]{1,5,5}}
  \proofstepwidestar[1]{f(e)\diamond f(e)=f(e)}{\rIE{9,8}}
  \proofstepwidestar[1]{f(e)\diamond u=f(e)}{%
    \FormulaRefAuto{GroupRightIdentity}[thm]{3,6}}
  \proofstepwidestar[1]{%
    f(e)\diamond u=f(e)\diamond f(e)}{\rIE{11,%
      \FormulaRefAuto{a=b\vdash b=a}{10}}}
  \proofstepwidestar[1]{u\in H}{%
      \FormulaRefAuto{GroupIdentityCarrier}[thm]{3}}
  \proofstepwidestar[1]{u=f(e)}{%
    \FormulaRefAuto{GroupLeftCancellation}[thm]{3,6,
      13,6,12}}
  \proofstepwidestar[1]{f(e)=u}{%
    \FormulaRefAuto{a=b\vdash b=a}{14}}
\end{tabproofwide}

\FormulaThmDeltaK[Gruppenhomomorphismen erhalten Inverse]{%
  \GroupHom{f}{G,\circ,e}{H,\diamond,u}\dsep x\in G
  \vdash f(x^{-1})=f(x)^{-1}%
}{GroupHomInversePreservation}{%
  \DeltaRow{Mengen}{G\dsep H\dsep e\dsep u\dsep x}
  \DeltaRow{Funktionen}{f\dsep\circ\dsep\diamond}
}
\begingroup
\ProofWideReasonsNearFormula
\begin{tabproofsplitwide}
  \proofpartwide{Träger der Bilder}
  \proofstepwidestarbreakkeep[1]{\GroupHom{f}{G,\circ,e}{H,\diamond,u}}{%
    \rA}
  \proofstepwidestarbreakkeep[2]{x\in G}{%
    \rA}
  \proofstepwidestarbreakkeep[1]{\GroupAlg{G}{\circ}{e}}{%
    \FormulaRefAuto{GroupHomDef}[def]{1}}
  \proofstepwidestarbreakkeep[1]{\GroupAlg{H}{\diamond}{u}}{%
    \FormulaRefAuto{GroupHomDef}[def]{1}}
  \proofstepwidestarbreakkeep[1]{f\colon G\to H}{%
    \FormulaRefAuto{GroupHomFunctionAxiom}[ax]{1}}
  \proofstepwidestarbreakkeep[1,2]{x^{-1}\in G}{%
    \FormulaRefAuto{GroupInverseCarrier}[thm]{3,2}}
  \proofstepwidestarbreakkeep[1,2]{f(x)\in H}{%
    \FormulaRefAuto{F\colon A\to B\dsep a\in A\vdash F(a)\in B}[thm]{5,2}}
  \proofstepwidestarbreakkeep[1,2]{f(x^{-1})\in H}{%
    \FormulaRefAuto{F\colon A\to B\dsep a\in A\vdash F(a)\in B}[thm]{5,6}}
  \closeproofpartwide

  \proofpartwide{Linksinverse und Eindeutigkeit}
  \setcounter{proofstepnr}{8}
  \proofstepwidestarbreakkeep[1]{f(e)=u}{%
    \FormulaRefAuto{GroupHomIdentityPreservation}[thm]{1}}
  \proofstepwidestarbreakkeep[1,2]{x^{-1}\circ x=e}{%
    \FormulaRefAuto{GroupInverseLeft}[thm]{3,2}}
  \proofstepwidestarbreakkeep[1,2]{f(x^{-1}\circ x)=f(x^{-1})\diamond f(x)}{%
    \FormulaRefAuto{GroupHomOperationAxiom}[ax]{1,6,2}}
  \proofstepwidestarbreakkeep[1,2]{f(x^{-1}\circ x)=u}{%
    \rIE{10,9}}
  \proofstepwidestarbreakkeep[1,2]{f(x^{-1})\diamond f(x)=u}{%
    \FormulaRefAuto{a=b\dsep a=c\vdash b=c}[thm]{11,12}}
  \proofstepwidestarbreakkeep[1,2]{f(x)^{-1}\in H}{%
    \FormulaRefAuto{GroupInverseCarrier}[thm]{4,7}}
  \proofstepwidestarbreakkeep[1,2]{f(x)\diamond f(x)^{-1}=u}{%
    \FormulaRefAuto{GroupInverseRight}[thm]{4,7}}
  \proofstepwidestarbreakkeep[1,2]{f(x^{-1})=f(x)^{-1}}{%
    \FormulaRefAuto{GroupInverseUnique}[thm]{4,7,8,14,13,15}}
  \closeproofpartwide
\end{tabproofsplitwide}
\endgroup

\FormulaThmDeltaK[Ein Gruppenhomomorphismus ist ein Monoidhomomorphismus]{%
  \GroupHom{f}{G,\circ,e}{H,\diamond,u}
  \vdash\MonoidHom{f}{G,\circ,e}{H,\diamond,u}%
}{GroupHomIsMonoidHom}{%
  \DeltaRow{Mengen}{G\dsep H\dsep e\dsep u}
  \DeltaRow{Funktionen}{f\dsep\circ\dsep\diamond}
}
\begin{tabproof}
  \proofstep{1}{\GroupHom{f}{G,\circ,e}{H,\diamond,u}}{\rA}
  \proofstep{1}{\GroupAlg{G}{\circ}{e}}{%
    \FormulaRefAuto{GroupHomDef}[def]{1}}
  \proofstep{1}{\GroupAlg{H}{\diamond}{u}}{%
    \FormulaRefAuto{GroupHomDef}[def]{1}}
  \proofstep{1}{\MonoidAlg{G}{\circ}{e}}{%
    \FormulaRefAuto{GroupBaseMonoidAxiom}[ax]{2}}
  \proofstep{1}{\MonoidAlg{H}{\diamond}{u}}{%
    \FormulaRefAuto{GroupBaseMonoidAxiom}[ax]{3}}
  \proofstep{1}{\SemigroupHom{f}{G,\circ}{H,\diamond}}{%
    \FormulaRefAuto{GroupHomDef}[def]{1}}
  \proofstep{1}{f(e)=u}{%
    \FormulaRefAuto{GroupHomIdentityPreservation}[thm]{1}}
  \proofstep{1}{\MonoidHom{f}{G,\circ,e}{H,\diamond,u}}{%
    \FormulaRefAuto{MonoidHomDef}[def]{4,5,6,7}}
\end{tabproof}

\section{Gruppenisomorphismen und Komposition}

\FormulaDefDeltaK[Begriff des Gruppenisomorphismus]{%
  \GroupIso{f}{G,\circ,e}{H,\diamond,u}%
}{GroupIsoDef}{%
  \DeltaRow{Mengen}{G\dsep H\dsep e\dsep u}
  \DeltaRow{Funktionen}{f\dsep\circ\dsep\diamond}
  \DeltaRow{\textbf{Axiome}}{}
  \DeltaPrem{Gruppenhomomorphismus}{%
    \GroupHom{f}{G,\circ,e}{H,\diamond,u}}
  \DeltaPrem{Bijektivität}{f\colon G\bij H}
  \DeltaRow{\textbf{Neues Symbol}}{}
  \DeltaPrem{Gruppenisomorphismus}{%
    \GroupIso{f}{G,\circ,e}{H,\diamond,u}}
}

\FormulaThmDeltaK[Komposition von Gruppenhomomorphismen]{%
  \begin{aligned}[t]
    &\GroupHom{f}{G,\circ,e}{H,\diamond,u}\dsep
      \GroupHom{g}{H,\diamond,u}{K,\bullet,v}\\
    &\vdash\GroupHom{g\circ f}{G,\circ,e}{K,\bullet,v}
  \end{aligned}
}{GroupHomComposition}{%
  \DeltaRow{Mengen}{G\dsep H\dsep K\dsep e\dsep u\dsep v}
  \DeltaRow{Funktionen}{f\dsep g\dsep\circ\dsep\diamond\dsep\bullet}
}
\begin{tabproof}
  \proofstep{1}{\GroupHom{f}{G,\circ,e}{H,\diamond,u}}{\rA}
  \proofstep{2}{\GroupHom{g}{H,\diamond,u}{K,\bullet,v}}{\rA}
  \proofstep{1}{\SemigroupHom{f}{G,\circ}{H,\diamond}}{%
    \FormulaRefAuto{GroupHomDef}[def]{1}}
  \proofstep{2}{\SemigroupHom{g}{H,\diamond}{K,\bullet}}{%
    \FormulaRefAuto{GroupHomDef}[def]{2}}
  \proofstep{1,2}{\SemigroupHom{g\circ f}{G,\circ}{K,\bullet}}{%
    \FormulaRefAuto{SemigroupHomComposition}[thm]{3,4}}
  \proofstep{1}{\GroupAlg{G}{\circ}{e}}{%
    \FormulaRefAuto{GroupHomDef}[def]{1}}
  \proofstep{2}{\GroupAlg{K}{\bullet}{v}}{%
    \FormulaRefAuto{GroupHomDef}[def]{2}}
  \proofstep{1,2}{\GroupHom{g\circ f}{G,\circ,e}{K,\bullet,v}}{%
    \FormulaRefAuto{GroupHomDef}[def]{6,7,5}}
\end{tabproof}

\FormulaThmDeltaK[Die Identität ist ein Gruppenisomorphismus]{%
  \GroupAlg{G}{\circ}{e}
  \vdash\GroupIso{\Id_G}{G,\circ,e}{G,\circ,e}%
}{GroupIdentityIsomorphism}{%
  \DeltaRow{Mengen}{G\dsep e}
  \DeltaRow{Funktionen}{\circ}
}
\begin{tabproof}
  \proofstep{1}{\GroupAlg{G}{\circ}{e}}{\rA}
  \proofstep{1}{\SemigroupAlg{G}{\circ}}{%
    \FormulaRefAuto{GroupBaseSemigroup}[thm]{1}}
  \proofstep{1}{\SemigroupIso{\Id_G}{G,\circ}{G,\circ}}{%
    \FormulaRefAuto{SemigroupIdentityIsomorphism}[thm]{2}}
  \proofstep{1}{\SemigroupHom{\Id_G}{G,\circ}{G,\circ}}{%
    \FormulaRefAuto{SemigroupIsoHomAxiom}[ax]{3}}
  \proofstep{1}{\GroupHom{\Id_G}{G,\circ,e}{G,\circ,e}}{%
    \FormulaRefAuto{GroupHomDef}[def]{1,1,4}}
  \proofstep{}{\Id_G\colon G\bij G}{%
    \FormulaRefAuto{\Id_A\colon A\bij A}[thm]}
  \proofstep{1}{\GroupIso{\Id_G}{G,\circ,e}{G,\circ,e}}{%
    \FormulaRefAuto{GroupIsoDef}[def]{5,6}}
\end{tabproof}

\section{Rekonstruktion aus der großen Potenzhalbgruppe}

\hypertarget{reconstruction.main}{}
Die \ReconstructionReadingLink{} erklärt zunächst die Erkennung der
Einermengen und anschließend, weshalb schon eine bekannte Ausgangsgruppe
für die Rekonstruktion genügt. Die zwölf zugehörigen formalen Ableitungen
stehen in den \ReconstructionProofsLink{}.
Ihre Aussagen, Nummern und bisherigen Referenzziele bleiben hier erhalten;
ebenso die allgemeinen Grundlagen zur Einheitengruppe eines Monoids.

\subsection{Der konkrete Transport von Gruppeneinermengen}

Für zwei Gruppen halten wir fest, welche Einermenge ein gegebener
Potenzisomorphismus wohin schickt. Diese Aussage wird anschließend
für die einseitige Gruppenstarrheit und später für Faserargumente verwendet.

\FormulaThmDeltaK[Jedes Gruppenelement ist eine Einheit]{%
  \GroupAlg{G}{\star}{e}\dsep g\in G\vdash\MonoidUnit{G,\star,e}{g}
}{GroupElementIsMonoidUnit}{%
  \DeltaRow{Mengen}{G\dsep e\dsep g\dsep h}
  \DeltaRow{Funktionen}{\star}
}
\begin{tabproof}
  \proofstep{1}{\GroupAlg{G}{\star}{e}}{\rA}
  \proofstep{2}{g\in G}{\rA}
  \proofstep{1}{\MonoidAlg{G}{\star}{e}}{\FormulaRefAuto{GroupBaseMonoidAxiom}[ax]{1}}
  \proofstep{1,2}{\exists h\in G\,(g\star h=e\land h\star g=e)}{\FormulaRefAuto{GroupInverseExistenceAxiom}[ax]{1,2}}
  \proofstep{1,2}{\MonoidUnit{G,\star,e}{g}}{\FormulaRefAuto{MonoidUnitDef}[def]{3,2,4}}
\end{tabproof}

\ReconstructionStatement{PowerGroupUnitsAreSingletons}
\par\noindent\ReconstructionProofLink{PowerGroupUnitsAreSingletons}.\par

\ReconstructionStatement{PowerGroupsSingletonLayer}
\par\noindent\ReconstructionProofLink{PowerGroupsSingletonLayer}.\par

\ReconstructionStatement{PowerGroupsSingletonTransport}
\par\noindent\ReconstructionProofLink{PowerGroupsSingletonTransport}.\par
\noindent\ReconstructionReadingLink{} mit vollständigem Lesebeweis.\par



\subsection{Einheitenmengen und Einheitenstabilität}

\FormulaDefDeltaK[Einheitenmenge eines Monoids]{%
  \MonoidAlg{A}{\star}{e}\vdash A^\times\coloneqq\{u\in A\mid\MonoidUnit{A,\star,e}{u}\}
}{MonoidUnitSetDef}{%
  \DeltaRow{Mengen}{A\dsep e\dsep u}
  \DeltaRow{Funktionen}{\star}
}

\FormulaThmDeltaK[Inverse bleiben in der Einheitenmenge]{%
  \MonoidAlg{A}{\star}{e}\dsep u\in A^\times\vdash\exists v\in A^\times\,(u\star v=e\land v\star u=e)
}{MonoidUnitSetInverse}{%
  \DeltaRow{Mengen}{A\dsep e\dsep u\dsep v}
  \DeltaRow{Funktionen}{\star}
}
\begingroup
\ProofWideReasonsNearFormula
\begin{tabproofwide}
  \proofstepwidestarbreakkeep[1]{\MonoidAlg{A}{\star}{e}}{%
    \rA}
  \proofstepwidestarbreakkeep[2]{u\in A^\times}{%
    \rA}
  \proofstepwidestarbreakkeep[1,2]{u\in A}{%
    \FormulaRefAuto{MonoidUnitSetDef}[def]{1,2}}
  \proofstepwidestarbreakkeep[1,2]{\MonoidUnit{A,\star,e}{u}}{%
    \FormulaRefAuto{MonoidUnitSetDef}[def]{1,2}}
  \proofstepwidestarbreakkeep[1,2]{\exists v\in A\,(u\star v=e\land v\star u=e)}{%
    \FormulaRefAuto{MonoidUnitDef}[def]{4}}
  \proofstepwidestarbreakkeep[6]{v\in A\land(u\star v=e\land v\star u=e)}{%
    \rA}
  \proofstepwidestarbreakkeep[6]{v\in A}{%
    \rAEa{6}}
  \proofstepwidestarbreakkeep[6]{u\star v=e\land v\star u=e}{%
    \rAEb{6}}
  \proofstepwidestarbreakkeep[1,2,6]{\MonoidUnit{A,\star,e}{v}}{%
    \FormulaRefAuto{MonoidUnitDef}[def]{1,7,\rEI{\rAI{3,\rAI{\rAEb{8},\rAEa{8}}}}}}
  \proofstepwidestarbreakkeep[1,2,6]{v\in A^\times}{%
    \FormulaRefAuto{MonoidUnitSetDef}[def]{1,7,9}}
  \proofstepwidestarbreakkeep[1,2,6]{\exists v\in A^\times\,(u\star v=e\land v\star u=e)}{%
    \rEI{\rAI{10,8}}}
  \proofstepwidestarbreakkeep[1,2]{\exists v\in A^\times\,(u\star v=e\land v\star u=e)}{%
    \rEE{5,6,11}}
\end{tabproofwide}
\endgroup

\FormulaThmDeltaK[Inverse eines Produkts im Monoid]{%
  \begin{aligned}[t]&\MonoidAlg{A}{\star}{e}\dsep u,v,u',v'\in A\dsep{}\\&u\star u'=e\dsep u'\star u=e\dsep v\star v'=e\dsep v'\star v=e\vdash{}\\&(u\star v)\star(v'\star u')=e\land(v'\star u')\star(u\star v)=e\end{aligned}
}{MonoidInversePairProduct}{%
  \DeltaRow{Mengen}{A\dsep e\dsep u\dsep v\dsep u'\dsep v'}
  \DeltaRow{Funktionen}{\star}
}
\begingroup
\ProofWideReasonsNearFormula
\begin{tabproofsplitwide}
  \proofpartwide{Voraussetzungen}
  \proofstepwidestarbreakkeep[1]{\MonoidAlg{A}{\star}{e}}{%
    \rA}
  \proofstepwidestarbreakkeep[2]{u\in A}{%
    \rA}
  \proofstepwidestarbreakkeep[3]{v\in A}{%
    \rA}
  \proofstepwidestarbreakkeep[4]{u'\in A}{%
    \rA}
  \proofstepwidestarbreakkeep[5]{v'\in A}{%
    \rA}
  \proofstepwidestarbreakkeep[6]{u\star u'=e}{%
    \rA}
  \proofstepwidestarbreakkeep[7]{u'\star u=e}{%
    \rA}
  \proofstepwidestarbreakkeep[8]{v\star v'=e}{%
    \rA}
  \proofstepwidestarbreakkeep[9]{v'\star v=e}{%
    \rA}
  \proofstepwidestarbreakkeep[1]{\SemigroupAlg{A}{\star}}{%
    \FormulaRefAuto{MonoidBaseSemigroupAxiom}[ax]{1}}
  \closeproofpartwide

  \proofpartwide{Rechtsinverse des Produkts}
  \setcounter{proofstepnr}{10}
  \proofstepwidestarbreakkeep[1,4,5]{v'\star u'\in A}{%
    \FormulaRefAuto{SemigroupClosureAxiom}[ax]{10,5,4}}
  \proofstepwidestarbreakkeep[1,2,3,4,5]{(u\star v)\star(v'\star u')=u\star(v\star(v'\star u'))}{%
    \FormulaRefAuto{SemigroupAssociativityAxiom}[ax]{10,2,3,11}}
  \proofstepwidestarbreakkeep[1,3,4,5]{(v\star v')\star u'=v\star(v'\star u')}{%
    \FormulaRefAuto{SemigroupAssociativityAxiom}[ax]{10,3,5,4}}
  \proofstepwidestarbreakkeep[1,3,4,5,8]{u\star(v\star(v'\star u'))=u\star(e\star u')}{%
    \rIE{13,8}}
  \proofstepwidestarbreakkeep[1,4]{e\star u'=u'}{%
    \FormulaRefAuto{MonoidLeftIdentityAxiom}[ax]{1,4}}
  \proofstepwidestarbreakkeep[1,4,6]{u\star(e\star u')=e}{%
    \rIE{15,6}}
  \proofstepwidestarbreakkeep[1,2,3,4,5,6,8]{(u\star v)\star(v'\star u')=e}{%
    \rChain{12,14,16}}
  \closeproofpartwide

  \proofpartwide{Linksinverse des Produkts}
  \setcounter{proofstepnr}{17}
  \proofstepwidestarbreakkeep[1,2,3]{u\star v\in A}{%
    \FormulaRefAuto{SemigroupClosureAxiom}[ax]{10,2,3}}
  \proofstepwidestarbreakkeep[1,2,3,4,5]{(v'\star u')\star(u\star v)=v'\star(u'\star(u\star v))}{%
    \FormulaRefAuto{SemigroupAssociativityAxiom}[ax]{10,5,4,18}}
  \proofstepwidestarbreakkeep[1,2,3,4]{(u'\star u)\star v=u'\star(u\star v)}{%
    \FormulaRefAuto{SemigroupAssociativityAxiom}[ax]{10,4,2,3}}
  \proofstepwidestarbreakkeep[1,2,3,4,7]{v'\star(u'\star(u\star v))=v'\star(e\star v)}{%
    \rIE{20,7}}
  \proofstepwidestarbreakkeep[1,3]{e\star v=v}{%
    \FormulaRefAuto{MonoidLeftIdentityAxiom}[ax]{1,3}}
  \proofstepwidestarbreakkeep[1,3,9]{v'\star(e\star v)=e}{%
    \rIE{22,9}}
  \proofstepwidestarbreakkeep[1,2,3,4,5,7,9]{(v'\star u')\star(u\star v)=e}{%
    \rChain{19,21,23}}
  \proofstepwidestarbreakkeep[1,2,3,4,5,6,7,8,9]{(u\star v)\star(v'\star u')=e\land(v'\star u')\star(u\star v)=e}{%
    \rAI{17,24}}
  \closeproofpartwide
\end{tabproofsplitwide}
\endgroup

\FormulaThmDeltaK[Produktabschluss der Einheitenmenge]{%
  \MonoidAlg{A}{\star}{e}\dsep u,v\in A^\times\vdash u\star v\in A^\times
}{MonoidUnitSetProduct}{%
  \DeltaRow{Mengen}{A\dsep e\dsep u\dsep v\dsep u'\dsep v'}
  \DeltaRow{Funktionen}{\star}
}
\begingroup
\ProofWideReasonsNearFormula
\begin{tabproofsplitwide}
  \proofpartwide{Inverse der Faktoren}
  \proofstepwidestarbreakkeep[1]{\MonoidAlg{A}{\star}{e}}{%
    \rA}
  \proofstepwidestarbreakkeep[2]{u\in A^\times}{%
    \rA}
  \proofstepwidestarbreakkeep[3]{v\in A^\times}{%
    \rA}
  \proofstepwidestarbreakkeep[1,2]{u\in A}{%
    \FormulaRefAuto{MonoidUnitSetDef}[def]{1,2}}
  \proofstepwidestarbreakkeep[1,2]{\exists u'\in A^\times\,(u\star u'=e\land u'\star u=e)}{%
    \FormulaRefAuto{MonoidUnitSetInverse}[thm]{1,2}}
  \proofstepwidestarbreakkeep[1,3]{v\in A}{%
    \FormulaRefAuto{MonoidUnitSetDef}[def]{1,3}}
  \proofstepwidestarbreakkeep[1,3]{\exists v'\in A^\times\,(v\star v'=e\land v'\star v=e)}{%
    \FormulaRefAuto{MonoidUnitSetInverse}[thm]{1,3}}
  \proofstepwidestarbreakkeep[8]{u'\in A^\times\land(u\star u'=e\land u'\star u=e)}{%
    \rA}
  \proofstepwidestarbreakkeep[9]{v'\in A^\times\land(v\star v'=e\land v'\star v=e)}{%
    \rA}
  \proofstepwidestarbreakkeep[1,8]{u'\in A}{%
    \FormulaRefAuto{MonoidUnitSetDef}[def]{1,\rAEa{8}}}
  \proofstepwidestarbreakkeep[1,9]{v'\in A}{%
    \FormulaRefAuto{MonoidUnitSetDef}[def]{1,\rAEa{9}}}
  \closeproofpartwide

  \proofpartwide{Produktinverse und Abschluss}
  \setcounter{proofstepnr}{11}
  \proofstepwidestarbreakkeep[1]{\SemigroupAlg{A}{\star}}{%
    \FormulaRefAuto{MonoidBaseSemigroupAxiom}[ax]{1}}
  \proofstepwidestarbreakkeep[1,2,3]{u\star v\in A}{%
    \FormulaRefAuto{SemigroupClosureAxiom}[ax]{12,4,6}}
  \proofstepwidestarbreakkeep[1,8,9]{v'\star u'\in A}{%
    \FormulaRefAuto{SemigroupClosureAxiom}[ax]{12,11,10}}
  \proofstepwidestarbreakkeep[1,2,3,8,9]{(u\star v)\star(v'\star u')=e\land(v'\star u')\star(u\star v)=e}{%
    \FormulaRefAuto{MonoidInversePairProduct}[thm]{1,4,6,10,11,\rAEa{\rAEb{8}},\rAEb{\rAEb{8}},\rAEa{\rAEb{9}},\rAEb{\rAEb{9}}}}
  \proofstepwidestarbreakkeep[1,2,3,8,9]{\MonoidUnit{A,\star,e}{u\star v}}{%
    \FormulaRefAuto{MonoidUnitDef}[def]{1,13,\rEI{\rAI{14,15}}}}
  \proofstepwidestarbreakkeep[1,2,3,8,9]{u\star v\in A^\times}{%
    \FormulaRefAuto{MonoidUnitSetDef}[def]{1,13,16}}
  \proofstepwidestarbreakkeep[1,2,3,8]{u\star v\in A^\times}{%
    \rEE{7,9,17}}
  \proofstepwidestarbreakkeep[1,2,3]{u\star v\in A^\times}{%
    \rEE{5,8,18}}
  \closeproofpartwide
\end{tabproofsplitwide}
\endgroup

\FormulaThmDeltaK[Träger und Nichtleerheit der Einheitenmenge]{%
  \MonoidAlg{A}{\star}{e}\vdash e\in A^\times\land A^\times\in\PowNE A
}{MonoidUnitSetCarrier}{%
  \DeltaRow{Mengen}{A\dsep e\dsep x\dsep y\dsep z}
  \DeltaRow{Funktionen}{\star}
}
\begingroup
\ProofWideReasonsNearFormula
\begin{tabproofwide}
  \proofstepwidestarbreakkeep[1]{\MonoidAlg{A}{\star}{e}}{%
    \rA}
  \proofstepwidestarbreakkeep[1]{e\in A}{%
    \FormulaRefAuto{MonoidIdentityCarrierAxiom}[ax]{1}}
  \proofstepwidestarbreakkeep[1]{\MonoidUnit{A,\star,e}{e}}{%
    \FormulaRefAuto{MonoidIdentityIsUnit}[thm]{1}}
  \proofstepwidestarbreakkeep[1]{e\in A^\times}{%
    \FormulaRefAuto{MonoidUnitSetDef}[def]{1,2,3}}
  \proofstepwidestarbreakkeep[5]{x\in A^\times}{%
    \rA}
  \proofstepwidestarbreakkeep[1,5]{x\in A}{%
    \FormulaRefAuto{MonoidUnitSetDef}[def]{1,5}}
  \proofstepwidestarbreakkeep[1]{A^\times\subseteq A}{%
    \FormulaRefAuto{SubsetIntroductionRule}[thm]{[5]\vdots 6}}
  \proofstepwidestarbreakkeep[1]{A^\times\neq\varnothing}{%
    \FormulaRefAuto{a\in A\vdash A\neq\varnothing}[thm]{4}}
  \proofstepwidestarbreakkeep[1]{A^\times\in\PowNE A}{%
    \FormulaRefAuto{NonemptyPowersetMembership}[thm]{\rAI{7,8}}}
  \proofstepwidestarbreakkeep[1]{e\in A^\times\land A^\times\in\PowNE A}{%
    \rAI{4,9}}
\end{tabproofwide}
\endgroup

\FormulaThmDeltaK[Die Einheiten bilden eine Gruppe]{%
  \MonoidAlg{A}{\star}{e}\vdash\GroupAlg{A^\times}{\star}{e}
}{MonoidUnitSetIsGroup}{%
  \DeltaRow{Mengen}{A\dsep e\dsep x\dsep y\dsep z}
  \DeltaRow{Funktionen}{\star}
}
\begingroup
\ProofWideReasonsNearFormula
\begin{tabproofsplitwide}
  \proofpartwide{Halbgruppe der Einheiten}
  \proofstepwidestarbreakkeep[1]{\MonoidAlg{A}{\star}{e}}{%
    \rA}
  \proofstepwidestarbreakkeep[1]{\SemigroupAlg{A}{\star}}{%
    \FormulaRefAuto{MonoidBaseSemigroupAxiom}[ax]{1}}
  \proofstepwidestarbreakkeep[1]{e\in A^\times\land A^\times\in\PowNE A}{%
    \FormulaRefAuto{MonoidUnitSetCarrier}[thm]{1}}
  \proofstepwidestarbreakkeep[1]{e\in A^\times}{%
    \rAEa{3}}
  \proofstepwidestarbreakkeep[1]{A^\times\subseteq A\land A^\times\neq\varnothing}{%
    \FormulaRefAuto{NonemptyPowersetMembership}[thm]{\rAEb{3}}}
  \proofstepwidestarbreakkeep[6]{x\in A^\times}{%
    \rA}
  \proofstepwidestarbreakkeep[7]{y\in A^\times}{%
    \rA}
  \proofstepwidestarbreakkeep[1,6,7]{x\star y\in A^\times}{%
    \FormulaRefAuto{MonoidUnitSetProduct}[thm]{1,6,7}}
  \proofstepwidestarbreakkeep[1]{\forall x\in A^\times\,\forall y\in A^\times\,(x\star y\in A^\times)}{%
    \rUI{\rRI{6,\rUI{\rRI{7,8}}}}}
  \proofstepwidestarbreakkeep[1]{\mathsf{UHGr}(A^\times;A,\star)}{%
    \FormulaRefAuto{SemigroupSubsemigroupCriterion}[thm]{2,\rAI{\rAEa{5},\rAI{\rAEb{5},9}}}}
  \proofstepwidestarbreakkeep[1]{\SemigroupAlg{A^\times}{\star}}{%
    \FormulaRefAuto{SubsemigroupFormsSemigroup}[thm]{2,10}}
  \closeproofpartwide

  \proofpartwide{Neutrales Element und Inverse}
  \setcounter{proofstepnr}{11}
  \proofstepwidestarbreakkeep[12]{z\in A^\times}{%
    \rA}
  \proofstepwidestarbreakkeep[1,12]{z\in A}{%
    \FormulaRefAuto{MonoidUnitSetDef}[def]{1,12}}
  \proofstepwidestarbreakkeep[1,12]{e\star z=z}{%
    \FormulaRefAuto{MonoidLeftIdentityAxiom}[ax]{1,13}}
  \proofstepwidestarbreakkeep[1,12]{z\star e=z}{%
    \FormulaRefAuto{MonoidRightIdentityAxiom}[ax]{1,13}}
  \proofstepwidestarbreakkeep[1]{\forall z\in A^\times\,(e\star z=z\land z\star e=z)}{%
    \rUI{\rRI{12,\rAI{14,15}}}}
  \proofstepwidestarbreakkeep[1]{\MonoidAlg{A^\times}{\star}{e}}{%
    \FormulaRefAuto{MonoidDef}[def]{11,4,16}}
  \proofstepwidestarbreakkeep[18]{w\in A^\times}{%
    \rA}
  \proofstepwidestarbreakkeep[1,18]{\exists y\in A^\times\,(w\star y=e\land y\star w=e)}{%
    \FormulaRefAuto{MonoidUnitSetInverse}[thm]{1,18}}
  \proofstepwidestarbreakkeep[1]{\forall w\in A^\times\,\exists y\in A^\times\,(w\star y=e\land y\star w=e)}{%
    \rUI{\rRI{18,19}}}
  \proofstepwidestarbreakkeep[1]{\GroupAlg{A^\times}{\star}{e}}{%
    \FormulaRefAuto{GroupDef}[def]{17,20}}
  \closeproofpartwide
\end{tabproofsplitwide}
\endgroup

\ReconstructionStatement{GroupPowerFullCarrierAbsorption}
\par\noindent\ReconstructionProofLink{GroupPowerFullCarrierAbsorption}.\par

\ReconstructionProductMemberStatement

\FormulaDefDeltaK[Einheitenstabiles Element]{%
  \begin{aligned}[t]&\operatorname{UStab}_{M,\star,e}(z)\coloneqq{}\\&z\in M\land\forall u\in M\,\bigl(\MonoidUnit{M,\star,e}{u}\rightarrow(u\star z=z\land z\star u=z)\bigr)\end{aligned}
}{MonoidUnitStableDef}{%
  \DeltaRow{Mengen}{M\dsep e\dsep z\dsep u}
  \DeltaRow{Funktionen}{\star}
}

\ReconstructionStatement{SemigroupIsoUnitStabilityTransport}
\par\noindent\ReconstructionProofLink{SemigroupIsoUnitStabilityTransport}.\par

\ReconstructionStatement{PowerMonoidUnitSetStable}
\par\noindent\ReconstructionProofLink{PowerMonoidUnitSetStable}.\par

\ReconstructionStatement{PowerMonoidUnitStableAbsorption}
\par\noindent\ReconstructionProofLink{PowerMonoidUnitStableAbsorption}.\par

\ReconstructionStatement{PowerMonoidUnitSubsetCriterion}
\par\noindent\ReconstructionProofLink{PowerMonoidUnitSubsetCriterion}.\par

\ReconstructionStatement{LargePowerUnitSetImage}
\par\noindent\ReconstructionProofLink{LargePowerUnitSetImage}.\par

\begingroup\small
\ReconstructionStatement{LargePowerUnitSubsetImage}
\endgroup
\par\noindent\ReconstructionProofLink{LargePowerUnitSubsetImage}.\par

\FormulaThmDeltaK[Alle Gruppenelemente sind Einheiten]{%
  \GroupAlg{A}{\star}{e}\vdash A^\times=A
}{GroupUnitSetIsCarrier}{%
  \DeltaRow{Mengen}{A\dsep e\dsep x\dsep y\dsep z}
  \DeltaRow{Funktionen}{\star}
}
\begingroup
\ProofWideReasonsNearFormula
\begin{tabproofwide}
  \proofstepwidestarbreakkeep[1]{\GroupAlg{A}{\star}{e}}{%
    \rA}
  \proofstepwidestarbreakkeep[1]{\MonoidAlg{A}{\star}{e}}{%
    \FormulaRefAuto{GroupBaseMonoidAxiom}[ax]{1}}
  \proofstepwidestarbreakkeep[1]{e\in A^\times\land A^\times\in\PowNE A}{%
    \FormulaRefAuto{MonoidUnitSetCarrier}[thm]{2}}
  \proofstepwidestarbreakkeep[1]{A^\times\subseteq A}{%
    \rAEa{\FormulaRefAuto{NonemptyPowersetMembership}[thm]{\rAEb{3}}}}
  \proofstepwidestarbreakkeep[5]{x\in A}{%
    \rA}
  \proofstepwidestarbreakkeep[1,5]{\MonoidUnit{A,\star,e}{x}}{%
    \FormulaRefAuto{GroupElementIsMonoidUnit}[thm]{1,5}}
  \proofstepwidestarbreakkeep[1,5]{x\in A^\times}{%
    \FormulaRefAuto{MonoidUnitSetDef}[def]{2,5,6}}
  \proofstepwidestarbreakkeep[1]{A\subseteq A^\times}{%
    \FormulaRefAuto{SubsetIntroductionRule}[thm]{[5]\vdots 7}}
  \proofstepwidestarbreakkeep[1]{A^\times=A}{%
    \FormulaRefAuto{A\subseteq B\land B\subseteq A\eqvdash A=B}[thm]{\rAI{4,8}}}
\end{tabproofwide}
\endgroup

\ReconstructionStatement{LargePowerSemigroupUnitGroupReconstruction}
\par\noindent\ReconstructionProofLink{LargePowerSemigroupUnitGroupReconstruction}.\par


Der folgende Gruppenstarrheitssatz geht in dieser Form auf
\href{https://doi.org/10.48550/arXiv.2606.01917}%
  {Liu--Tringali (2026), Satz~1.3}
zurück. Er setzt nur auf einer Seite eine Gruppe voraus; auf der anderen
Seite genügt zunächst eine Halbgruppe. Eine Endlichkeitsannahme ist nicht nötig.

\ReconstructionStatement{GroupPowerSemigroupRigidity}
\par\noindent\ReconstructionProofLink{GroupPowerSemigroupRigidity}.\par
\noindent\ReconstructionReadingLink{} mit vollständigem Lesebeweis.\par


\begin{remark}
Der Beweis verwendet wesentlich die \emph{große} Potenzhalbgruppe aller
nichtleeren Teilmengen: Die ganze Einheitengruppe \(A^\times\) muss selbst
als Element von \(\PowNE A\) auftreten.  Für endliche Grundträger stimmt
die große mit der finitären Potenzhalbgruppe überein.  Bei unendlichen
Gruppen folgt die entsprechende Aussage für die Halbgruppe nur der
endlichen nichtleeren Teilmengen dagegen nicht aus diesem Beweis.
\end{remark}

\chapter{Ergänzende Isomorphiesätze}

\section{Umkehrung und Transport von Gruppenisomorphismen}

Die folgenden Sätze halten zwei allgemeine Werkzeuge ausdrücklich fest:
Gruppenisomorphismen lassen sich umkehren, und die Gruppenstruktur wird
schon durch einen Halbgruppenisomorphismus auf den Zielträger übertragen.
Im zweiten Fall muss die Zielstruktur vorher noch nicht als Gruppe bekannt
sein.

\FormulaThmDeltaK[Die Umkehrabbildung eines Gruppenisomorphismus]{%
  \begin{aligned}[t]
    &\GroupIso{f}{A,\star,e}{B,\diamond,u}\vdash{}\\[-2pt]
    &\GroupIso{f^{-1}}{B,\diamond,u}{A,\star,e}
  \end{aligned}
}{GroupIsoInverse}{%
  \DeltaRow{Mengen}{A\dsep B\dsep e\dsep u}
  \DeltaRow{Funktionen}{f\dsep f^{-1}\dsep\star\dsep\diamond}
}
\begingroup
\ProofWideReasonsNearFormula
\setlength{\LTpre}{0pt}
\begin{tabproofwide}
  \proofstepwidestarbreakkeep[1]{\GroupIso{f}{A,\star,e}{B,\diamond,u}}{\rA}
  \proofstepwidestarbreakkeep[1]{\GroupHom{f}{A,\star,e}{B,\diamond,u}}{%
    \FormulaRefAuto{GroupIsoDef}[def]{1}}
  \proofstepwidestarbreakkeep[1]{\GroupAlg{A}{\star}{e}}{%
    \FormulaRefAuto{GroupHomDef}[def]{2}}
  \proofstepwidestarbreakkeep[1]{\GroupAlg{B}{\diamond}{u}}{%
    \FormulaRefAuto{GroupHomDef}[def]{2}}
  \proofstepwidestarbreakkeep[1]{\SemigroupHom{f}{A,\star}{B,\diamond}}{%
    \FormulaRefAuto{GroupHomDef}[def]{2}}
  \proofstepwidestarbreakkeep[1]{f\colon A\bij B}{%
    \FormulaRefAuto{GroupIsoDef}[def]{1}}
  \proofstepwidestarbreakkeep[1]{\SemigroupIso{f}{A,\star}{B,\diamond}}{%
    \FormulaRefAuto{SemigroupIsoDef}[def]{5,6}}
  \proofstepwidestarbreakkeep[1]{\SemigroupIso{f^{-1}}{B,\diamond}{A,\star}}{%
    \FormulaRefAuto{SemigroupIsoInverse}[thm]{7}}
  \proofstepwidestarbreakkeep[1]{\SemigroupHom{f^{-1}}{B,\diamond}{A,\star}}{%
    \FormulaRefAuto{SemigroupIsoHomAxiom}[ax]{8}}
  \proofstepwidestarbreakkeep[1]{f^{-1}\colon B\bij A}{%
    \FormulaRefAuto{SemigroupIsoBijectiveAxiom}[ax]{8}}
  \proofstepwidestarbreakkeep[1]{\GroupHom{f^{-1}}{B,\diamond,u}{A,\star,e}}{%
    \FormulaRefAuto{GroupHomDef}[def]{4,3,9}}
  \proofstepwidestarbreakkeep[1]{\GroupIso{f^{-1}}{B,\diamond,u}{A,\star,e}}{%
    \FormulaRefAuto{GroupIsoDef}[def]{11,10}}
\end{tabproofwide}
\endgroup

\FormulaThmDeltaK[Ein Halbgruppenisomorphismus transportiert die Gruppenstruktur]{%
  \begin{aligned}[t]
    &\GroupAlg{A}{\star}{e}\dsep
      \SemigroupIso{f}{A,\star}{B,\diamond}\vdash{}\\[-2pt]
    &\GroupAlg{B}{\diamond}{f(e)}
  \end{aligned}
}{SemigroupIsoGroupTransport}{%
  \DeltaRow{Mengen}{A\dsep B\dsep e\dsep a\dsep c\dsep y\dsep z}
  \DeltaRow{Funktionen}{f\dsep\star\dsep\diamond}
}
\begingroup
\ProofWideReasonsNearFormula
\begin{tabproofsplitwide}
  \proofpartwide{Monoidstruktur und Urbilder}
  \proofstepwidestarbreakkeep[1]{\GroupAlg{A}{\star}{e}}{\rA}
  \proofstepwidestarbreakkeep[2]{\SemigroupIso{f}{A,\star}{B,\diamond}}{\rA}
  \proofstepwidestarbreakkeep[1]{\MonoidAlg{A}{\star}{e}}{%
    \FormulaRefAuto{GroupBaseMonoidAxiom}[ax]{1}}
  \proofstepwidestarbreakkeep[1,2]{\MonoidAlg{B}{\diamond}{f(e)}}{%
    \FormulaRefAuto{SemigroupIsoMonoidTransport}[thm]{3,2}}
  \proofstepwidestarbreakkeep[2]{f\colon A\bij B}{%
    \FormulaRefAuto{SemigroupIsoBijectiveAxiom}[ax]{2}}
  \proofstepwidestarbreakkeep[6]{y\in B}{\rA}
  \proofstepwidestarbreakkeep[2,6]{\exists a\in A\,(f(a)=y)}{%
    \FormulaRefAuto{Bijektive Funktion}[def]{5,6}}
  \proofstepwidestarbreakkeep[8]{a\in A\land f(a)=y}{\rA}
  \proofstepwidestarbreakkeep[1,8]{\exists c\in A\,(a\star c=e\land c\star a=e)}{%
    \FormulaRefAuto{GroupInverseExistenceAxiom}[ax]{1,\rAEa{8}}}
  \closeproofpartwide

  \proofpartwide{Inverse im Ziel und Gruppenschluss}
  \setcounter{proofstepnr}{9}
  \proofstepwidestarbreakkeep[10]{c\in A\land(a\star c=e\land c\star a=e)}{\rA}
  \proofstepwidestarbreakkeep[2,10]{f(c)\in B}{%
    \FormulaRefAuto{F\colon A\bij B\dsep x\in A\vdash F(x)\in B}{5,\rAEa{10}}}
  \proofstepwidestarbreakkeep[1,2,8,10]{%
    \begin{aligned}[t]
      &(a\star c=e\land c\star a=e)\leftrightarrow{}\\[-2pt]
      &(f(a)\diamond f(c)=f(e)\land f(c)\diamond f(a)=f(e))
    \end{aligned}}{%
    \FormulaRefAuto{SemigroupIsoMonoidInversePair}[thm]{3,4,2,\rAEa{8},\rAEa{10}}}
  \proofstepwidestarbreakkeep[1,2,8,10]{%
    f(a)\diamond f(c)=f(e)\land f(c)\diamond f(a)=f(e)}{%
    \rRE{\rLREa{12},\rAEb{10}}}
  \proofstepwidestarbreakkeep[1,2,8,10]{%
    y\diamond f(c)=f(e)\land f(c)\diamond y=f(e)}{%
    \rIE{\rAEb{8},13}}
  \proofstepwidestarbreakkeep[1,2,8,10]{%
    \exists z\in B\,(y\diamond z=f(e)\land z\diamond y=f(e))}{%
    \rEI{\rAI{11,14}}}
  \proofstepwidestarbreakkeep[1,2,8]{%
    \exists z\in B\,(y\diamond z=f(e)\land z\diamond y=f(e))}{%
    \rEE{9,10,15}}
  \proofstepwidestarbreakkeep[1,2,6]{%
    \exists z\in B\,(y\diamond z=f(e)\land z\diamond y=f(e))}{%
    \rEE{7,8,16}}
  \proofstepwidestarbreakkeep[1,2]{%
    \forall y\in B\,\exists z\in B\,
      (y\diamond z=f(e)\land z\diamond y=f(e))}{%
    \rUI{\rRI{6,17}}}
  \proofstepwidestarbreakkeep[1,2]{\GroupAlg{B}{\diamond}{f(e)}}{%
    \FormulaRefAuto{GroupDef}[def]{4,18}}
  \closeproofpartwide
\end{tabproofsplitwide}
\endgroup

\end{document}
